%%%%%%%%%%%%%%%%%%%%%%%%%%%%%%%%%%%%%%%%%
% Important note:
% This template must be compiled with XeLaTeX, the below lines will ensure this
%!TEX TS-program = xelatex
%!TEX encoding = UTF-8 Unicode
%
%%%%%%%%%%%%%%%%%%%%%%%%%%%%%%%%%%%%%%%%%

%----------------------------------------------------------------------------------------
%	PACKAGES AND OTHER DOCUMENT CONFIGURATIONS
%----------------------------------------------------------------------------------------

\documentclass[11pt, letterpaper]{awesome-cv} % A4 paper size by default, use 'letterpaper' for US letter

\geometry{left=2cm, top=1.5cm, right=2cm, bottom=2cm, footskip=.5cm} % Configure page margins with geometry
 
\fontdir[fonts/] % Specify the location of the included fonts

% Color for highlights
%\colorlet{awesome}{awesome-skyblue} % Default colors include: awesome-emerald, awesome-skyblue, awesome-red, awesome-pink, awesome-orange, awesome-nephritis, awesome-concrete, awesome-darknight
%\definecolor{awesome}{HTML}{CA63A8} % Uncomment if you would like to specify your own color
\definecolor{awesome}{HTML}{0BA89A} % UCSB sea green

% Colors for text - uncomment and modify
%\definecolor{darktext}{HTML}{000000}
\definecolor{text}{HTML}{000000}
\definecolor{graytext}{HTML}{000000}
%\definecolor{lighttext}{HTML}{000000}

\renewcommand{\acvHeaderSocialSep}{\quad\textbar\quad} % If you would like to change the social information separator from a pipe (|) to something else

%----------------------------------------------------------------------------------------
%	PERSONAL INFORMATION
%	Comment any of the lines below if they are not required
%----------------------------------------------------------------------------------------

\name{Jason C.}{Wei}
\address{Pleasanton, California, USA}
\mobile{(925) 352-7600}
\github{Jayweiii}
\homepage{jweii.com}
\email{jwei@ucsb.edu}
\linkedin{jwei}
%\skype{skypeid}
%\stackoverflow{SOid}{SOname}
%\twitter{@twit}
%\linkedin{linkedin name}
%\reddit{reddit account}
%\xing{xing name}
%\extrainfo{test} % Other text you want to include on this line

\position{Mechanical Engineer{\enskip\cdotp\enskip}UCSB} % Your expertise/fields
\quote{Second-year MechE at UCSB with extensive hands-on experience in \\ 3D printing, Formula SAE EV powertrain design, phase-change thermofluidics, and CNC Milling.} % A quote or statement

%----------------------------------------------------------------------------------------
%	RECIPIENT/POSITION/LETTER INFORMATION
%	All of the below lines must be filled out
%----------------------------------------------------------------------------------------

\recipient{}{Begley Research Lab\\Santa Barbara, CA} % The company being applied to

\letterdate{\today} % The date on the letter, default is the date of compilation

\lettertitle{Undergraduate Researcher} % The title of the letter

\letteropening{Dear Begley Research Lab,} % How the letter is opened

\letterclosing{Kind regards,} % How the letter is closed

% \letterenclosure[Attached]{Resume} % Any enclosures with the letter

\makecvfooter{\today}{Jason C. Wei~~~·~~~Cover Letter}{} % Specify the letter footer with 3 arguments: (<left>, <center>, <right>), leave any of these blank if they are not needed
  
%----------------------------------------------------------------------------------------

\begin{document}

\makecvheader % Print the header

\makelettertitle % Print the title

%----------------------------------------------------------------------------------------
%	LETTER CONTENT
%----------------------------------------------------------------------------------------

\begin{cvletter}

%------------------------------------------------

I am delighted to submit my application for the 3D Printing Undergraduate Research opportunity at the Begley Research Lab. As a second-year Mechanical Engineering student, my extensive background in 3D printing, FEA, and thermofluidic micro-structures makes me a strong candidate to contribute to your group's research in structural micro-structures.

In my role as an undergraduate researcher at Yangying's Thermofluid Energy Research Laboratory, I have been contributing to their research of analyzing flow instabilities due to micro-structures. I have worked very closely with PhD student Harsh Shah and manufactured the entire testing apparatus using 3D printing and CNC milling. Additionally, as a personal research project under the mentorship of Yangying, I have been exploring combining 3D printed internal geometries, such as gyroids, with gas-liquid double phase change, to increase the efficiency of heat exchangers. 

As a Workshop Wizard and 3D Printer Tool Lead at the CNSI Microfluidics lab and Innovation Workshop, I manage, repair, and improve processes for our Haas Super Mini Mill and all types of our 3D printers, including FDM, SLA, DLP, and PolyJet. My hands-on experiences in manufacturing researchers' experimental setups, including Harsh's setup that I described above, has honed my understanding of the different design and material considerations of additive manufacturing, subtractive manufacturing (CNC \& manual machining), chemical post-processes (such as acetone vapor smoothing), and even combining these different manufacturing methods in creating a single part. I am also familiar with our Raman spectroscopy system, as I have been working on an improved light shielding cover to make it easier to use the machine.

At Gaucho Racing, I lead the design and manufacturing of our safe, modular, and repairable high voltage battery pack, drivetrain, and mounting solutions. To optimize our structures for minimal mass, I use Ansys FEA and 3D print with carbon fiber-filled nylon using my personal 3D printer. I have explored with using different infill and support geometries, including AI generated ones, for different strength and cost applications.

I am excited to bring the skills that I have accrued through hands-on engineering experiences to the Begley Research Lab. Through this research opportunity, I hope to meaningfully contribute to the field of additive manufacturing and making structures stronger, lighter, and less expensive. Thank you for considering my application. I look forward to discussing this research opportunity and my qualifications in further detail.

%------------------------------------------------

\end{cvletter}

%----------------------------------------------------------------------------------------

\makeletterclosing % Print the signature and enclosures

\end{document}